\documentclass[czech, kiv, sem, he, iso690alph, pdf,viewonly]{fasthesis}
\title{\texttt{Virtuální herbář funkcí}}
\author{Marek}{Jirušek}{}{}
\supervisor{Ing. Petr Vaněček, Ph.D.}
\nobastardtitle
\nocopyrightnotice
\addbibresource{references.bib}
\begin{document}

\frontpages[notm]
\tableofcontents


\chapter{Úvod}
Pochopit chování funkcí o jedné proměnné je pro studenty středních škol relativně intuitivní, jelikož lze snadno vizualizovat jejich grafy na papíře nebo pomocí, v dnešní době již jednoduše dostupných, nástrojů.
Chování funkcí o dvou a více proměnných však již není tak jednoduché, jelikož nejenže vyžaduje určitou představivost, ale v dnešní době neexistuje dostatečný počet dostupných nástrojů, které by umožňovali interaktivní zkoumání těchto funkcí.
Tato práce vznikla zejména z důvodu digitalizovat herbář funkcí, který byl napsán na katedře matematiky.
Cílem této práce bylo vytvořit virtuální herbář funkcí, který umožní studentům interaktivně prozkoumávat chování funkcí o dvou proměnných ve virtuální realitě.

\section{Prohlášení o použití umělé inteligence}
Během přípravy této práce autor použil nástroj umělé inteligence Google Gemini k následujícím účelům:
\begin{itemize}
    \item konzultace týkající se implementačních detailů a návrhu softwarové architektury systému,
    \item asistence při formulaci a stylistické úpravě odborného textu a kontrola terminologické konzistence.
\end{itemize}
Po použití tohoto nástroje autor pečlivě zkontroloval a upravil veškerý obsah podle potřeby. 
Veškerá rozhodnutí týkající se návrhu, implementace a interpretace výsledků jsou výhradně vlastní prací autora. Autor přebírá plnou odpovědnost za výsledný obsah práce.

\section{Definice problému}
Generování 3D modelů grafů funkcí v reálném čase je náročný úkol, jelikož vyžaduje spoustu výpočtů.
Navíc, pro zajištění plynulého zážitku ve virtuální realitě, je důležitý výkon a minimální odezva při překreslování grafů po změně parametrů.
Aplikace musí být schopna:
\begin{itemize}
    \item Zpracovat uživatelský vstup: Aplikace musí být schopna přijímat nové funkce, zpracovávat změny parametrů a definičního oboru a aktualizovat grafy v reálném čase.
    \item Generování 3D modelů: Aplikace musí být schopna generovat 3D modely grafů funkcí, které jsou vizuálně přehledné a informativní a zároveň edukativní pro studenty.
    \item Body nespojitosti: Aplikace by měla identifikovat a vizualizovat body nespojitosti funkcí.
    \item Interakce ve VR: aplikace by měla umožnit uživatelům za běhu měnit předpisy funkcí, měnit parametry a definiční obor a okamžitě vidět aktualizované grafy ve virtuální realitě bez opouštění VR prostředí.
\end{itemize}

\section{Cíl práce}
Tento bakalářský projekt ideově a obsahově vychází z publikace \textit{Herbář funkcí}, která vznikla na Katedře matematiky (FAV/KMA). 
Hlavním cílem práce je posunout tento koncept do interaktivního trojrozměrného prostoru a vyvinout aplikaci „Virtuální herbář funkcí“, 
která uživatelům umožní prozkoumávat matematické funkce dvou proměnných v prostředí virtuální reality (VR). 

Aplikace je navržena jako názorná učební pomůcka, jež má studentům usnadnit pochopení prostorových grafů a demonstrovat, 
jak se jejich tvar dynamicky mění v závislosti na různých parametrech a zvoleném definičním oboru.

Pro úspěšné dosažení tohoto hlavního cíle byly stanoveny následující dílčí úkoly:
\begin{itemize}
    \item Implementace algoritmu pro převod textového zápisu funkce do podoby 3D sítě (meshe) a její plynulé generování v reálném čase.
    \item Návrh řešení pro detekci bodů nespojitosti a korektní ořezávání asymptot, aby nedocházelo k deformacím generovaného modelu.
    \item Optimalizace výpočetní náročnosti aplikace s cílem udržet stabilní snímkovou frekvenci nezbytnou pro VR.
    \item Vytvoření prostorového uživatelského rozhraní umožňujícího intuitivní manipulaci s modelem, úpravu parametrů rovnice a změnu definičního oboru přímo ve virtuálním prostředí.
\end{itemize}

\section{Struktura dokumentace}
Následující text je logicky členěn do několika celků, které popisují vývoj a implementaci aplikace. 
Druhá kapitola představuje celkovou architekturu aplikace a způsob, jakým mezi sebou komunikují jednotlivé komponenty. 
Třetí kapitola poté rozebírá matematické jádro, od zpracování textového vstupu přes kompilaci do abstraktního syntaktického stromu (AST) až po samotný algoritmus generování 3D sítě a řešení asymptot. 
Čtvrtá kapitola se následně věnuje optimalizačním technikám a limitům enginu, které bylo nutné překonat pro zajištění plynulého chodu ve VR. 
Pátá kapitola popisuje návrh prostorového uživatelského rozhraní a implementaci fyzikálních interakcí i volného pohybu pomocí rámce XR Interaction Toolkit.
Poslední kapitoly se zabývají testováním a shrnutím celého projektu, včetně návrhů pro budoucí rozšíření a vylepšení.


\chapter{Architektura aplikace}
Architektura aplikace je strukturována do specializovaných modulů, které spolu komunikují primárně pomocí událostí (C\# Events) a centrálního zprostředkovatele. 
Návrh dodržuje princip jedné odpovědnosti (\textit{Single Responsibility Principle}). Tento přístup zajišťuje snadnou udržitelnost kódu, zabraňuje vzniku takzvaných 
\textit{God Objects} a umožňuje plynulý chod aplikace i při výpočetně náročných operacích ve virtuální realitě.

\section{Komponenty aplikace}
Aplikace je logicky rozdělena do několika hlavních softwarových komponent, z nichž každá plní svou specifickou roli:

\begin{itemize}
    \item \textbf{Manager:} Představuje ústřední řídicí uzel aplikace (využívající návrhový vzor \textit{Mediator}). 
    Sám neprovádí složité výpočty, ale drží reference na ostatní subsystémy a koordinuje jejich součinnost.
    \item \textbf{DataManager:} Zajišťuje vrstvu perzistence. 
    Stará se o načítání textového souboru s rovnicemi ze souborového systému a na pozadí monitoruje jeho případné externí změny.
    \item \textbf{VRInputManager:} Zapouzdřuje logiku ovládání ve virtuální realitě. Izoluje čtení vstupů z ovladačů a spravuje fyzikální interakce uživatele s modelem, jako je například obouruční škálování velikosti grafu (přes sub-komponentu \command{HandInteractionHandler}).
    \item \textbf{Generator:} Tvoří matematické jádro aplikace (fungující jako vzor \textit{Facade}). 
    Přijímá textové řetězce reprezentující funkce a deleguje jejich zpracování -- kompilaci provádí \command{MathPreprocess}, iterativní výpočet prostorových souřadnic a detekci asymptot řeší 
    \command{MeshGenerator} a stavbu os \command{KSSGenerator}.
    \item \textbf{UIManager a MenuGenerator:} Starají se o dynamické prostorové rozhraní. 
    Moduly naslouchají událostem z \command{Generatoru} a \command{DataManageru} a na jejich základě automaticky instancují ovládací prvky (informační tabule, posuvníky a číselníky) 
    pro manipulaci s definičním oborem a parametry.
\end{itemize}

\section{Datový tok a komunikace}
Aby nedocházelo k propadům snímkové frekvence, není komunikace mezi moduly vázána na pravidelnou aktualizační smyčku enginu, ale je striktně řízena událostmi (event-driven architecture). 

Zpracování nové funkce probíhá asynchronně. 
Jakmile \command{DataManager} detekuje nebo načte funkce, vyvolá událost \command{OnFunctionsLoaded}. 
Na tu zareaguje \command{UIManager} překreslením nabídky na virtuální tabuli a \command{Generator} zahájením kompilace první funkce. 
Po úspěšném sestavení 3D sítě vyvolá \command{Generator} událost \command{OnFunctionLoaded}. Tuto událost zachytí \command{MenuGenerator} a 
dynamicky podle ní sestaví příslušné ovládací prvky pro zachycené parametry rovnice.

Pokud následně uživatel pomocí paprsku ovladače změní hodnotu parametru, menu zapíše novou hodnotu přímo do pole parametrů v \command{Generatoru} a explicitně zavolá metodu \command{RefreshGraph()}. 
Aplikace tak negeneruje síť neustále, ale výhradně v okamžiku, kdy dojde ke skutečné změně vstupních dat.

\chapter{Architektura dat a matematické jádro}
Tato kapitola se věnuje vnitřní logice aplikace, která zajišťuje převod textových řetězců s matematickými výrazy do podoby interaktivních 3D modelů ve virtuální realitě. 
Architektura je datově řízená (data-driven), tím pádem veškeré předpisy funkcí jsou uloženy externě a aplikace je dynamicky zpracovává až za běhu. 
Toto řešení umožňuje uživatelům průběžně modifikovat obsah virtuálního herbáře bez nutnosti zásahu do zdrojového kódu či restartování samotné aplikace, což usnadňuje proces výuky a experimentování.

\section{Datově řízený návrh a real-time monitoring}
Základním architektonickým principem řešení je striktní oddělení výpočetní logiky od spravovaných dat. 
Místo statického definování tzv. \textit{hard-coded} funkcí přímo v modulu \command{Generator} využívá aplikace externí textový konfigurační soubor s názvem \filename{functions.txt}, 
ve kterém je definován seznam rovnic tvořících obsah herbáře.

\subsection{Úložiště a správa souborů}
Centrálním úložištěm pro veškeré matematické definice je soubor \filename{functions.txt}, který je alokován ve speciálním adresáři \filename{StreamingAssets}. 
Tato složka je z pohledu enginu \product{Unity} specifická tím, že její obsah nepodléhá interní kompilaci a zůstává v nezměněné podobě i po sestavení výsledné aplikace (buildu). 
Koncoví uživatelé, typicky vyučující, tak mohou snadno doplňovat či upravovat vyučovanou látku zcela nezávisle na vývojovém cyklu programu.

Při startu aplikace provádí třída \command{DataManager} kontrolu integrity datového úložiště. Pokud aplikace detekuje, že konfigurační soubor chybí nebo je poškozen, 
automaticky vygeneruje nový a zainicializuje jej sadou předdefinovaných demonstračních rovnic. 
Pro maximální zjednodušení uživatelského přístupu navíc aplikace při prvním spuštění vytvoří v kořenovém adresáři programu zástupce ve formě dávkového souboru \filename{1\_Upravit\_Rovnice.bat}, 
který zajistí okamžité otevření konfiguračního souboru ve výchozím textovém editoru hostitelského operačního systému.

\subsection{Mechanismus FileSystemWatcher a asynchronní aktualizace}
Aby bylo možné modifikovat rovnice za běhu aplikace, využívá \command{DataManager} zapouzdřenou .NET komponentu \command{FileSystemWatcher} (prostřednictvím vlastní třídy \command{FileWatcher}). 
Tento nástroj na pozadí monitoruje cílový adresář a naslouchá události zápisu (\command{NotifyFilters.LastWrite}) do souboru \filename{functions.txt}.

Jelikož události ze sledování souborového systému probíhají ve vedlejším výpočetním vlákně a aplikační rozhraní enginu \product{Unity} není takzvaně \textit{thread-safe} (bezpečné pro přímý přístup z více vláken), 
je proces aktualizace rozdělen do dvou synchronizovaných fází:

\begin{itemize}
    \item \textbf{Fáze detekce (Vlákno na pozadí):} Jakmile \command{FileSystemWatcher} zaznamená změnu souboru, je spuštěna anonymní metoda, která pouze změní hodnotu booleovské proměnné \command{needsRefresh} na hodnotu \command{true}. Během této fáze se nevolají žádné metody enginu Unity, čímž je zamezeno kolizím a pádům aplikace.
    \item \textbf{Fáze zpracování (Hlavní vlákno Unity):} V rámci standardní herní smyčky (metoda \command{Update}) \command{DataManager} každý snímek kontroluje stav proměnné \command{needsRefresh}. Pokud je indikována změna, proměnná je bezpečně resetována a následně je z hlavního vlákna zavolána metoda pro opětovné načtení souboru a aktualizaci aplikace.
\end{itemize}

Tento \textit{Thread-Safe} mechanismus zajišťuje maximální stabilitu aplikace i při rychlém a opakovaném přepisování konfiguračního souboru externími programy.

\section{Analýza matematických výrazů a kompilace do AST}
Jakmile je textová reprezentace funkce (např. výraz $a \cdot \sin(b \cdot x)$) načtena, je předána třídě \command{MathPreprocess}. 
Ta zajišťuje její převod do reverzní polské notace (RPN) a následnou kompilaci do abstraktního syntaktického stromu (AST). 

Vzhledem k tomu, že vstupem je textový řetězec, nasazení prostého interpretu by bylo pro potřeby aplikace zcela nedostačující. 
Opakované vyhodnocování textu pro desítky tisíc bodů sítě v každém vykreslovaném snímku by vedlo k enormní výpočetní režii a drastickému propadu snímkové frekvence, 
což je ve virtuální realitě kritický problém způsobující kinetózu (nevolnost). Zpracování výrazu je proto rozděleno do následujících tří fází:

\begin{itemize}
    \item \textbf{Lexikální analýza (Tokenizace):} Rozdělení vstupního řetězce na elementární textové a sémantické jednotky (tzv. tokeny).
    \item \textbf{Shunting-yard algoritmus:} Syntaktická analýza tokenů a jejich převod z infixové notace do RPN, čímž je jednoznačně určena priorita operací bez nutnosti použití závorek.
    \item \textbf{Kompilace do AST:} Konstrukce syntaktického stromu z RPN a jeho překlad do nativního kódu (do podoby delegátu) pro dosažení maximální rychlosti při samotném výpočtu bodů.
\end{itemize}

\subsection{Lexikální analýza (Tokenizace)}
Než může být matematický výraz algoritmicky zpracován, musí být rozložen na základní stavební bloky, tzv. tokeny. Úkolem metody \command{Tokenize} je lineárně projít řetězec a rozdělit jej na jednotky, 
se kterými lze následně pracovat.

Před samotným průchodem je text očištěn od bílých znaků (mezer), převeden na malá písmena a případné desetinné čárky jsou normalizovány na tečky, čímž se zajistí kompatibilita s 
parsováním datového typu \command{float} v jazyce C\#. Následně algoritmus iteruje řetězcem znak po znaku a shlukuje je do následujících tří kategorií:

\begin{itemize}
    \item \textbf{Numerické hodnoty:} Narazí-li analyzátor na číslici nebo desetinnou tečku, shlukuje všechny bezprostředně následující platné znaky do jediného číselného tokenu.
    \item \textbf{Alfanumerické řetězce (Slova):} Detekuje-li analyzátor písmeno, spojí souvislou řadu písmen do jednoho textového tokenu. V případě výrazu $a \cdot \sin(b \cdot x)$ 
    tak budou rozpoznány samostatné textové jednotky \command{a}, \command{sin}, \command{b} a \command{x}.
    \item \textbf{Operátory a závorky:} Jednotlivé znaky z předem definované množiny \command{+-*/^()} jsou izolovány jako samostatné matematické či strukturní operátory.
\end{itemize}

Kromě základního rozdělování textu se analyzátor v této fázi elegantně vypořádává s problematickým unárním mínusem, který se v matematických výrazech běžně vyskytuje 
(např. ve výrazech \command{-x}, \command{-(x+y)} nebo \command{sin(-x)}). Analyzátor tuto situaci uvnitř cyklu řeší zavedením jednoduchého pravidla: pokud je detekován znak \command{-} 
na samém začátku výrazu nebo bezprostředně za otevírací závorkou \command{()}, je automaticky před tento znak vložen nový token \command{0}. 
Tím dochází k internímu převodu unární operace na operaci binární (ve smyslu $0 - x$), což efektivně a bezpečně předchází pádům při vyhodnocování priorit v navazujícím algoritmu Shunting-yard.

\subsection{Algoritmus Shunting-yard}
Výstupem předchozí fáze je seznam tokenů v běžné infixové notaci. Tato notace je přirozená pro člověka, pro přímé strojové vyhodnocení je však nevhodná z důvodu nutnosti explicitně řešit priority 
operátorů a uzávorkování. Pro převod výrazu do reverzní polské notace (RPN) implementuje třída \command{MathPreprocess} modifikovaný algoritmus \term{Shunting-yard} \cite{dijkstra_shunting_yard}.

V RPN se operátory zapisují za své operandy, čímž odpadá nutnost použití závorek. 
Metoda \command{ShuntingYard} pro tento převod využívá dvě datové struktury: výstupní seznam (\command{List<string> outputQueue}) a dočasný zásobník (\command{operatorStack}). 
Prioritu operátorů definuje pevný slovník; sčítání a odčítání mají nejnižší prioritu (1), umocňování nejvyšší (3).

Algoritmus zpracovává tokeny zleva doprava a operandy (čísla, proměnné) vkládá přímo do výstupu. Operátory a matematické funkce odkládá na dočasný zásobník, 
dokud není jednoznačně určena jejich priorita vůči dalším operacím. Při načtení matematického operátoru dojde k jeho porovnání s operátorem na vrcholu zásobníku. 
Pokud má operátor na zásobníku vyšší nebo stejnou prioritu, přesune se do výstupního seznamu. Výjimkou je operátor mocniny (\command{\^}), který je asociativní zprava doleva.

Závorky plní roli bariér. Otevírací závorka se ukládá na zásobník. Při nalezení zavírací závorky se všechny operátory na zásobníku až po první otevírací závorku přesunou na výstup. 
Po zpracování všech vstupních tokenů se zbylé operátory ze zásobníku vyprázdní do výstupu. Tím vzniká kompletní RPN reprezentace zadané funkce, připravená pro kompilaci do abstraktního syntaktického stromu (AST).

Oproti standardní verzi algoritmu, která zpracovává primárně základní aritmetické operátory a číselné konstanty, obsahuje použitá implementace dvě klíčové modifikace. 
První úpravou je podpora proměnných (například souřadnic $x$ a $z$). Ty lexikální analyzátor identifikuje a algoritmus je následně propouští přímo na výstup jako běžné operandy. 
Druhým rozšířením je nativní zpracování matematických funkcí (např. goniometrických nebo logaritmických). 
Tyto funkce se dočasně ukládají na zásobník a k jejich přesunu do výstupní fronty dochází cíleně po zpracování zavírací závorky argumentu. 
Tím je algoritmicky zaručeno správné zřetězení funkce s jejím vnitřním výrazem, což standardní Dijkstrův návrh nativně neřeší.

\begin{code}{C#}{Strukturální pseudokód algoritmu Shunting-yard\label{lst:shunting_yard}}
foreach (string token in vstupniTokeny) {
    if (JeOperand(token)) {
        outputQueue.Add(token);
    } else if (JeFunkce(token)) {
        operatorStack.Push(token);
    } else if (JeOperator(token)) {
        while (operatorStack.Count > 0 && operatorStack.Peek() != "(") {
            string topOp = operatorStack.Peek();
            if (token != "^" && Priorita(token) <= Priorita(topOp)) {
                outputQueue.Add(operatorStack.Pop());
            } else {
                break;
            }
        }
        operatorStack.Push(token);
    } else if (token == "(") {
        operatorStack.Push(token);
    } else if (token == ")") {
        while (operatorStack.Count > 0 && operatorStack.Peek() != "(") {
            outputQueue.Add(operatorStack.Pop());
        }
        if (operatorStack.Count > 0 && operatorStack.Peek() == "(") {
            operatorStack.Pop(); // Zahození levé závorky
        }
        if (operatorStack.Count > 0 && JeFunkce(operatorStack.Peek())) {
            outputQueue.Add(operatorStack.Pop());
        }
    }
}

while (operatorStack.Count > 0) {
    outputQueue.Add(operatorStack.Pop());
}
\end{code}

\subsection{Kompilace do abstraktního syntaktického stromu}
Závěrečnou fází předzpracování funkce je transformace fronty v reverzní polské notaci do podoby optimalizovaného spustitelného delegátu. 
Cílem této fáze je sestavit v paměti takovou strukturu, kterou bude procesor schopen opakovaně a maximálně efektivně vykonávat v reálném čase.

\subsubsection{Abstraktní syntaktický strom (AST)}
Abstraktní syntaktický strom (z anglického \term{Abstract Syntax Tree}, zkráceně AST) je hierarchická stromová struktura reprezentující logickou stavbu zdrojového kódu nebo matematického výrazu. 
Vnitřní uzly stromu tvoří operátory (například sčítání či násobení) a matematické funkce (například sinus nebo absolutní hodnota). Listy stromu představují operandy, tedy číselné konstanty nebo proměnné.

Výhodou AST je absence nadbytečných syntaktických prvků, jako jsou závorky, jelikož posloupnost vyhodnocování je dána samotnou hierarchií stromu. Výpočet probíhá rekurzivně od nejnižších listů směrem ke kořeni.

\subsubsection{Konstrukce stromu a kompilace v aplikaci}
Proces stavby AST probíhá lineárním průchodem RPN fronty za využití knihovny \product{System.Linq.Expressions} \cite{microsoft_docs_expressions} a pomocného zásobníku typu \command{Stack<Expression>}. 
Algoritmus postupně odebírá tokeny a podle jejich typu vytváří příslušné uzly:

\begin{itemize}
    \item \textbf{Konstanty a proměnné:} Číselné hodnoty jsou transformovány na uzly typu \command{ConstantExpression}. Prostorové proměnné $x$ a $z$ jsou definovány jako uzly \command{ParameterExpression}, sloužící jako argumenty pro iterativní výpočet.
    \item \textbf{Parametry (koeficienty):} Pokud analyzátor narazí na identifikátor, který nepatří mezi interní funkce (například koeficienty $a, b$), přiřadí mu index v dynamickém poli parametrů. Přístup k těmto hodnotám je reprezentován uzlem \command{Expression.ArrayIndex} nad polem \command{argsParam}.
    \item \textbf{Aritmetické operátory:} Při detekci operátoru dojde k vyjmutí dvou předchozích výrazů ze zásobníku a jejich spojení do nového uzlu (například \command{Expression.Add}). Vzniklý podstrom je uložen zpět na zásobník.
    \item \textbf{Mapování funkcí a reflexe:} Vzhledem k tomu, že použitá knihovna neobsahuje nativní uzly pro goniometrické funkce nebo mocniny, využívá kompilátor mechanismus reflexe jazyka C\#. Pomocí metody \command{GetMethod} jsou vyhledány odpovídající definice z knihovny \product{UnityEngine.Mathf} (například \command{Sin}, \command{Pow}). Tyto operace jsou do stromu vloženy jako uzly typu \command{MethodCallExpression}.
\end{itemize}

Výsledkem procesu je kořenový uzel stromu zapouzdřený v objektu \command{LambdaExpression}. Voláním metody \command{.Compile()} dochází k JIT (Just-In-Time) kompilaci do mezijazyka \term{Intermediate Language} (IL). Tím vzniká spustitelný předpis typu \command{Func<float, float, float[], float>}. Metoda \command{Compile} vrací tento předpis společně se seznamem názvů nalezených koeficientů, který aplikace využívá pro dynamické generování ovládacích prvků v uživatelském rozhraní.


\subsubsection{Optimalizace zobrazení pomocí pomocné paměti (Caching)}
Nejprostším způsobem vyhodnocení matematického výrazu by bylo naprogramování \term{interpretu}, který by pro každý počítaný bod sítě znovu analyzoval a zpracovával vstupní textový řetězec. Vzhledem k tomu, že 3D síť grafu může běžně obsahovat i více než 100\,000 vrcholů a prostředí virtuální reality vyžaduje plynulé překreslování s frekvencí 72 až 90 Hz, musel by takový interpret provádět několik zbytečných textových operací každou vteřinu. To by vedlo k extrémní výpočetní režii a naprosté degradaci plynulosti obrazu.

Dynamická kompilace do IL kódu tuto zátěž prakticky zcela eliminuje. Zkompilovaný delegát je procesorem vykonáván nativně a s výkonem plně srovnatelným s kódem, který by programátor napevno zapsal do zdrojového souboru.

Aby se zamezilo kompilování již využívaných předpisů, využívá třída \command{Generator} vnitřní mezipaměť (\term{Caching}). 
Každý úspěšně zkompilovaný předpis je společně se seznamem nalezených koeficientů trvale uložen do hashovací tabulky \command{Dictionary}, konkrétně do privátní kolekce \command{compiledCache}. 
Pokud se uživatel k dané rovnici v herbáři později vrátí, aplikace ji okamžitě načte přímo z paměti v její přeložené podobě. 
Tento přístup garantuje rychlou odezvu uživatelského rozhraní i při rychlém přepínání mezi složitými matematickými modely na virtuální tabuli.

\subsection{Detekce nespojitostí a algoritmus bisekce}
Většina matematických funkcí má složité chování, jako jsou například body nespojitosti, mezery v definičním oboru nebo asymptoty (např. $f(x,y) = \frac{x^2 \cdot y^2}{x^3 + y^2}$). 
V místech takových nespojitostí vznikají při naivním generování sítě vizuální artefakty v podobě nepřirozeně natažených a strmých polygonů. 
Pro korektní zobrazení těchto jevů využívá aplikace vícefázový vyhodnocovací proces doplněný o numerickou metodu \term{bisekce} (půlení intervalu).

Generování sítě proto neprobíhá naslepo. Každý vypočítaný bod prochází validací. 
Pokud zkompilovaná funkce vrátí neplatnou hodnotu (\command{NaN} nebo nekonečno) nebo vypočítaná výška přesáhne povolený limit prostoru (\command{limitY}), označí algoritmus daný vrchol jako neplatný. 

Následně algoritmus prozkoumá čtyři bezprostřední sousedy tohoto bodu v mřížce. Pokud alespoň jeden ze sousedů leží v platné oblasti, znamená to, že se výpočet nachází na okraji asymptoty či trhliny a spouští se proces bisekce. Princip hledání přesné hrany spočívá v postupném půlení vzdálenosti mezi platným a neplatným vrcholem:

\begin{enumerate}
    \item Algoritmus určí středový bod (\command{mid}) úsečky spojující platný a neplatný bod.
    \item Tento středový bod je obratem vyhodnocen pomocí zkompilovaného delegátu.
    \item Pokud má středový bod platnou výšku a shodné znaménko s původním platným bodem, posouvá se platná hranice přímo do tohoto středu.
    \item V opačném případě (bod je stále neplatný nebo jeho výpočet odskočil do opačného nekonečna) se do středu posouvá neplatná hranice.
\end{enumerate}

Tento cyklus se standardně zopakuje desetkrát (v kódu definováno konstantou \command{BINARY\_SEARCH\_ITERATIONS}). 
Původně neplatný bod je následně ukotven přesně na tyto nově spočítané souřadnice.

\subsection{Triangulace, směrové vektory a výškové ořezání}
Ukotvení bodů na hranu asymptoty řeší pouze přesnou topologii vrcholů. Neřeší však samotný problém nechtěného propojování polygonů napříč touto trhlinou. 
Z toho důvodu si algoritmus během bisekce u každého označeného bodu ukládá směrové příznaky \command{dirX} a \command{dirZ}. 
Tyto vektory indikují, ze kterého směru byl původně neplatný bod přitažen ke svému platnému sousedovi. Pokud byl například bod přitažen vlevo (\command{dirX = -1}), 
znamená to, že fyzicky patří k levému břehu asymptoty. Pokud záchrana bodu nebyla vůbec možná, je směrovému vektoru přiřazena hodnota \command{INVALID\_DIRECTION\_MARKER}.

Ve fázi triangulace, kdy se vrcholy spojují do finálních trojúhelníků, generátor tyto směrové vektory kontroluje. Základní plošný segment sítě je tvořen čtveřicí vrcholů (čtvercem). Algoritmus zkoumá, zda některý z vrcholů nebyl přitažen ze směru, který míří pryč od středu aktuálně řešeného čtverce. Pokud ano, znamená to, že se daný čtverec snaží překlenout mezeru v grafu. Algoritmus v takovém případě jeho vykreslení přeskočí příkazem \command{continue}, čímž tyto defektní polygony vyřadí.

Navíc se nezávisle ověřuje, zda maximální výškový rozdíl mezi rohy čtverce nepřesahuje kritickou toleranci a zda funkce v samotném středu čtverce nevrací neplatnou hodnotu. 
Zjednodušená logika této topologické ochrany, zbavená takzvaných magických čísel a využívající zavedené konstanty, je znázorněna v následujícím pseudokódu:

\begin{code}{C#}{Logika zamezení propojování polygonů přes asymptoty\label{lst:triangulation}}
// Prochazeni mrizky a ziskani rohu aktualniho ctverce:
// DL (Dolni-Levy), DP (Dolni-Pravy), HL (Horni-Levy), HP (Horni-Pravy)

// 1. Zahozeni ctvercu, ktere obsahuji zcela neplatne body
if (DL.y == INVALID_HEIGHT_MARKER || DP.y == INVALID_HEIGHT_MARKER || 
    HL.y == INVALID_HEIGHT_MARKER || HP.y == INVALID_HEIGHT_MARKER) continue;

// 2. Kontrola smerovych vektoru z bisekce (-1 pro vlevo/dolu, 1 pro vpravo/nahoru)
// Pokud byl vrchol pritazen z opacne strany, nez lezi stred ctverce, propast je uvnitr.
// Nebo pokud nesel zachranit (INVALID_DIRECTION_MARKER), tak jej automaticky preskoci.
if (DL.dirX == -1 || DL.dirZ == -1 || DL.dirX == INVALID_DIRECTION_MARKER) continue; 
if (DP.dirX ==  1 || DP.dirZ == -1 || DP.dirX == INVALID_DIRECTION_MARKER) continue; 
if (HL.dirX == -1 || HL.dirZ ==  1 || HL.dirX == INVALID_DIRECTION_MARKER) continue; 
if (HP.dirX ==  1 || HP.dirZ ==  1 || HP.dirX == INVALID_DIRECTION_MARKER) continue; 

// 3. Kontrola extremni strmosti 
if ((MaxVyska(DL, DP, HL, HP) - MinVyska(DL, DP, HL, HP)) > maxDifference) continue;

// 4. Pojistna kontrola stredu ctverce (odhaleni velmi uzkych asymptot)
float stredY = VypoctiFunkciV(stredX, stredZ);
if (float.IsNaN(stredY) || stredY > limitY) continue;

VytvorTrojuhelnikyCtverce(DL, DP, HL, HP);
\end{code}

Ačkoliv je výše zmíněnými kroky zajištěno, že nevznikají polygony překračující trhliny, je nutné korektně zpracovat i samotné vrcholy směřující k nekonečnu. 
Tyto okrajové platné vrcholy jsou fyzicky oříznuty na maximální povolenou výšku (\command{limitY}) pomocí metody \command{Mathf.Clamp}. Aby na těchto oříznutých místech nevznikaly nepřirozené vizuální stropy, 
je výsledná síť doplněna o speciální grafický materiál (shader). Ten v kombinaci s fyzickým ořezáním zajišťuje, že se trojúhelníky blížící se k okraji limitu vykreslují postupně zprůhledněné nebo jsou zcela zahozeny. Vytváří se tak plynulý a přirozený dojem funkce mizející v nekonečnu.


\subsection{Kompletace modelu a fyzikální interakce}
Vrcholy, které se nacházejí hluboko v neplatné oblasti a nebylo možné je zachránit algoritmem bisekce, si dočasně ponechávají přiřazenou extrémní výškovou hodnotu 
(definovanou pomocí konstanty \command{INVALID\_HEIGHT\_MARKER}). Těsně před finálním odesláním struktury do paměti grafické karty jsou tyto neviditelné vrcholy přesunuty přesně do lokálního počátku 
souřadného systému (\command{Vector3.zero}). Jedná se o vrcholy, které nakonec nejsou součástí žádného vykreslovaného trojúhelníku.

Tato normalizace datových polí je kritická z hlediska enginu pro výpočet fyziky. Jakmile je trojrozměrná síť kompletně vygenerována, je nad ní automaticky sestavena kolizní obálka \command{BoxCollider}. 
Ta je nezbytná pro správné vyhodnocování interakcí ve virtuální realitě. Její prostorové rozměry (\command{mesh.bounds.size}) a střed (\command{mesh.bounds.center}) 
jsou dynamicky vypočteny na základě reálných prostorových extrémů viditelného grafu. 

Schování neplatných bodů do lokálního počátku zajišťuje, že se vypočtená obálka nebude uměle natahovat do nekonečna. Díky tomuto kroku odpovídá fyzikální těleso vizuální reprezentaci grafu. 
Uživatel tak může ve virtuální realitě výsledný matematický objekt fyzicky uchopit a plynule s ním manipulovat prostřednictvím vývojového rámce \product{XR Interaction Toolkit}.





\chapter{Optimalizace a limity enginu}
Vývoj aplikací pro virtuální realitu je specifický tím, že vyžaduje nejen vysoký výpočetní výkon, ale i efektivní správu paměti a vysokou optimalizaci. 
Zatímco u běžných desktopových aplikací je občasný pokles snímkové frekvence pouze estetickým defektem, ve virtuální realitě vede jakékoliv zpoždění při vykreslování (tzv. \term{frame drop}) k závažným problémům. 
Může způsobovat nevolnosti a výrazně tak zhoršit celkový uživatelský zážitek.

Jelikož aplikace generuje v reálném čase modely o velikosti přesahující 10\,000 vrcholů a tyto modely uživatel dynamicky mění, bylo nutné přistoupit k optimalizaci na několika úrovních. 
Optimalizace se dotkla samotné architektury kódu, použitých datových struktur i specifických funkcí enginu \product{Unity}.

\section{Indexování sítě}
Základem každého 3D objektu v enginu \product{Unity} je datová struktura \command"Mesh". 
Ta se skládá z pole vrcholů (\command"vertices") a pole indexů (\command"triangles"), jež tyto vrcholy spojují do trojúhelníkových ploch. 

Z výkonnostních důvodů využívá engine \product{Unity} ve výchozím nastavení 16bitový formát pro indexování vrcholů. 
Tento formát pojme maximálně $2^{16}$ unikátních hodnot, což striktně omezuje velikost jednoho trojrozměrného modelu na 65\,536 vrcholů. 
Pro vizualizaci některých matematických funkcí, zejména těch s vysokou frekvencí změn nebo s komplexními tvary, je však potřeba generovat sítě s mnohem vyšším rozlišením.

Aby bylo možné tento hardwarový limit překonat, aplikace modifikuje vnitřní nastavení geometrické sítě. 
Před samotným naplněním daty je formát indexů přepnut do 32bitového režimu:

\begin{code}{C}{Aktivace 32bitového indexování geometrie\label{lst:index_format}}
mesh.indexFormat = UnityEngine.Rendering.IndexFormat.UInt32;
\end{code}

Toto nastavení posouvá teoretický limit jednoho objektu na více než 4 miliardy vrcholů. 
Díky tomuto kroku není rozlišení vizualizované mřížky limitováno samotným enginem, ale pouze reálným výpočetním výkonem hardwaru. 
To umožňuje škálovatelnost kvality zobrazení podle možností dané grafické karty a aktuálních požadavků uživatele.

\section{Správa paměti a eliminace alokací}
Dalším kritickým místem (\term{bottleneck}) při dynamickém generování 3D modelů je správa operační paměti a činnost tzv. \term{Garbage Collectoru} (GC). 
Pokud by aplikace při každé změně parametru (např. při posunutí posuvníku o jeden milimetr) vytvořila v paměti zcela nový objekt \command{Mesh} a nová pole pro vrcholy a indexy, 
došlo by nejen k rychlému vyčerpání dostupné paměti, ale i k častým zásahům GC, které by způsobovaly závažné poklesy snímkové frekvence a zhoršovaly by plynulost zobrazení.

Aplikace proto recykluje již existující objekty. 
Objekt \command{Mesh} je v paměti alokován pouze jednou, a to při inicializaci. 
Při požadavku na překreslení grafu se nevytváří objekt nový, ale existující struktura je pouze vyčištěna voláním metody \command{mesh.Clear()} a následně jsou do ní nalita nová, přepsaná data. 

Obdobně je optimalizován i výpočet normálových vektorů. 
Namísto manuálního počítání kolmic pro každý trojúhelník aplikace využívá vestavěnou funkci \command{mesh.RecalculateNormals()}, která je optimalizována pro rychlé zpracování velkých množství dat 
a minimalizuje režii spojenou s výpočtem normál.

\section{Architektura řízená událostmi (Event-driven)}
Jak již bylo naznačeno v předchozích kapitolách, klíčovou vrstvou optimalizace je odstranění náročných výpočtů z hlavní renderovací smyčky (metody \command{Update}). 
Celá architektura je proto řízena událostmi (\textit{event-driven}). 

Zatímco vizuální systém ve VR běží neustále a plynule překresluje obraz pro displeje brýlí, matematické jádro ve výchozím stavu spí. 
Až ve chvíli, kdy uživatel fyzicky změní stav některého z ovládacích prvků rozhraní (například posune posuvníkem parametru), zachytí příslušný modul (např. \command{MenuGenerator} v rámci \command{UIManageru}) 
tuto událost. Tento modul následně přímo zaktualizuje potřebné hodnoty v poli parametrů a explicitně zavolá překreslovací metodu v \command{Generatoru} (\command{RefreshGraph()}). 

Tento přístup garantuje, že procesorový čas je spotřebováván na matematické operace výhradně v momentech, kdy dochází k reálné změně vstupních dat. 
Po zbytek času zůstává plný výpočetní výkon k dispozici pro potřeby enginu a fyziky virtuální reality, což je kritické pro udržení stabilní a vysoké snímkové frekvence.



\chapter{Uživatelské rozhraní, interakce a zobrazení ve VR}
Návrh uživatelského rozhraní ve virtuální realitě vyžaduje oproti tradičním desktopovým aplikacím zcela odlišný přístup. 
Uživatel nevyužívá standardní vstupní zařízení, jako jsou myš a klávesnice, ale spoléhá se výhradně na prostorovou manipulaci. 
Pro implementaci těchto interakčních mechanismů aplikace využívá vývojový rámec \product{XR Interaction Toolkit} (XRIT) \cite{unity_xr_manual}, 
který je oficiální součástí herního enginu \product{Unity}. Tento balíček poskytuje abstrakci pro mapování vstupů z různých hardwarových ovladačů (například Meta Quest, HTC Vive nebo Valve Index) 
a zajišťuje jejich následný převod na interakční události.

\section{Prostorové uživatelské rozhraní (Tabule a příruční menu)}
Uživatelské rozhraní aplikace je ukotveno přímo do prostoru 3D scény ve formě interaktivních panelů (tzv. \term{World Space Canvas}). Z hlediska softwarové logiky a přístupnosti se dělí na dva hlavní celky:

\begin{itemize}
    \item \textbf{Hlavní tabule (Výběr funkce):} Tento statický panel je umístěný ve virtuálním prostoru a slouží k výběru analyzované rovnice. 
    Poté, co \command{DataManager} načte předpisy funkcí z externího textového souboru, předá je komponentě \command{TableManager}. 
    Ta následně vyčistí herní desku a pomocí metody \command{RefreshBoardUI()} vygeneruje seznam nově načtených interaktivních tlačítek (přepínačů). 
    Uživatel na tlačítko namíří laserový paprsek a stisknutím volí konkrétní funkci, kterou chce vizualizovat.
    
    \item \textbf{Příruční menu (Parametry a definiční obor):} Pro operativní manipulaci s vygenerovaným grafem má uživatel k dispozici prostorové menu. 
    Zapínání a vypínání jednotlivých panelů (aby se zamezilo vizuálnímu přeplnění scény) spravuje skript \command{HandMenuHandler} na základě stisku tlačítek na VR ovladači. 
    Celé menu spadá pod zastřešující prvek \command{UIManager} a dělí se na dvě specifické sekce:
    \begin{itemize}
        \item \textit{Menu pro úpravu parametrů:} Obsahuje komponentu \command{ParametrMenuHandler}, která na základě dat z generátoru automaticky instancuje posuvníky (\term{Slidery}) 
        pro proměnné koeficienty (například $a, b, c$) detekované v aktuálně zvolené rovnici.
        \item \textit{Menu definičního oboru:} Obsahuje komponentu \command{DomainMenuHandler}, 
        jež generuje krokovací tlačítka (\term{Steppery}) pro přesnou úpravu definičního oboru na osách $x$ a $z$ ($X_{min}, X_{max}, Z_{min}, Z_{max}$). 
        K dispozici je navíc speciální nastavení \textit{Uniform} pro symetrické a rovnoměrné škálování definičního oboru do všech směrů současně.
    \end{itemize}
\end{itemize}

Pro interakci s těmito panely se používá komponenta \product{XR Ray Interactor}, která je připojena k virtuálním reprezentacím obou ovladačů. 
Tato komponenta vrhá do scény virtuální laserový paprsek (\term{Raycast}). 
Jakmile paprsek protne prvek uživatelského rozhraní, aplikace zaregistruje kolizi a umožní s daným prvkem manipulovat pomocí tlačítek na ovladači. 
Díky architektuře řízené událostmi je jakýkoliv vstup z uživatelského rozhraní (například pohyb posuvníku parametru) okamžitě zachycen, hodnoty jsou zapsány do generátoru a graf je plynule překreslen voláním metody 
\command{RefreshGraph()}.


\section{Fyzikální manipulace a volný pohyb}
Aby bylo možné zkoumat topologii vygenerovaných funkcí ze všech představitelných úhlů, podporuje aplikace přímý fyzický úchop grafu a volný pohyb uživatele v prostoru. 
Kolem každého vygenerovaného grafu je automaticky obalena komponenta \command{BoxCollider}. 
Její střed a rozměry jsou po každé kompilaci dynamicky přepočítávány na základě aktuálních prostorových hranic (\term{bounds}) vygenerované sítě. 
Na tento kolizní objekt je navázaná komponenta \command{XRGrabInteractable}, která tvoří přímé rozhraní mezi uživatelovou rukou a 3D sítí grafu.

Logika fyzické interakce funguje na následujících principech:

\begin{itemize}
    \item \textbf{Přímý úchop (Direct Interaction):} Pokud uživatel fyzicky přiblíží svůj virtuální ovladač ke kolizní obálce grafu a stiskne úchopové tlačítko (\term{Grip Button}), objekt se stane dočasným potomkem ruky. Uživatel tak může graf zvednout ze stojanu, libovolně s ním v prostoru rotovat nebo si jej přiblížit k očím.
    
    \item \textbf{Úchop na dálku (Ray Interaction):} Graf není nutné zvedat pouze z bezprostřední blízkosti. Je možné jej uchopit a přesouvat i na dálku pomocí integrovaného virtuálního laserového paprsku.
    
    \item \textbf{Manuální změna měřítka (Scaling):} Pokud uživatel stiskne aktivační tlačítka na obou ovladačích současně, přepne se specializovaná komponenta \command{HandInteractionHandler} do režimu obouručního škálování. V tomto stavu skript kontinuálně měří měnící se euklidovskou vzdálenost mezi levým a pravým ovladačem. Pomocí přirozeného gesta roztahování rukou od sebe nebo k sobě tak uživatel může měřítko celého 3D grafu plynule zvětšovat či zmenšovat.
    
    \item \textbf{Zohlednění transformací a reset pozice:} Jakmile uživatel objekt pustí, interakční rámec zruší vazbu na ovladač a graf zůstane v prostoru ukotven přesně v pozici a rotaci, ve které byl zanechán. Pro případ, že by uživatel graf nechtěně odhodil nebo by objekt propadl mimo herní zónu, obsahuje třída \command{VRInputManager} metodu \command{ResetGraph()}. Ta po stisku určeného tlačítka na ovladači navrátí graf do výchozí pozice, vynuluje jeho fyzikální rychlosti (\texttt{Rigidbody}) a synchronně zresetuje matematický definiční obor i hodnoty parametrů.
    
    \item \textbf{Prostupné prostředí:} Uživatel není vázán pouze na statické pozorování modelu zvenčí. Kolizní systém je nastaven tak, že detekuje úchopy ovladačů, ale neklade odpor fyzickému pohybu hráče. Uživatel tak může volně procházet i skrz samotný vnitřek matematického grafu.
\end{itemize}

Díky tomu, že algoritmus bisekce přesouvá veškeré vadné vrcholy mířící do nekonečna bezpečně do lokálního počátku (\command{Vector3.zero}), 
odpovídá velikost dynamicky upraveného kolizního tělesa \command{BoxCollider} reálné vizuální podobě grafu. Nedochází tak k matoucím a nechtěným kolizím v prázdném prostoru okolo asymptot.


\section{Vizuální reprezentace a orientace v prostoru}
Pro správné didaktické pochopení matematického modelu je nezbytné, aby uživatel neztratil orientaci v souřadnicovém systému ani při libovolném natočení grafu. 
Vizuální prezentace modelu je proto podpořena následujícími prvky:

\begin{itemize}
    \item \textbf{Souřadnicové osy (KSS):} Každý vygenerovaný graf je obalen vizualizací základních prostorových os ($x, y, z$). 
    Tyto osy (reprezentované barevnými 3D válci), včetně dynamických číselných ticků a popisků, jsou generovány statickou třídou \command{KSSGenerator} prostřednictvím metody \command{DrawKSS()}. 
    V hierarchii scény jsou ukotveny jako přímí potomci grafu pod prázdným objektem \texttt{KSS\_Root}. 
    K textové reprezentaci hodnot na osách je využívána vysoce ostrá a optimalizovaná komponenta \product{TextMeshPro}.
    
    \item \textbf{Oboustranné vykreslování:} Pro vykreslování trojúhelníkové sítě byl vytvořen vlastní grafický program (\term{shader}). 
    Ten primárně zajišťuje oboustranné vykreslování polygonů (\term{Double-sided rendering}). 
    Tato vlastnost je zcela nezbytná pro zachování plné viditelnosti povrchu grafu i při zkoumání a procházení modelu z jeho vnitřní nebo spodní strany.
    
    \item \textbf{Dithering a optimalizovaná průhlednost:} Aby okraje grafu a oříznuté asymptoty nevypadaly jako nepřirozeně useknuté stěny, využívá materiál techniku \term{ditheringu} 
    (šumového rozptylu pixelů na základě matice). Při běžném poloprůhledném prolínání (\term{alpha blending}) dochází ve VR k vizuálnímu rozmazání, chybám v řazení průhledných ploch a ztrátě prostorových detailů. 
    Dithering tento problém eliminuje tím, že některé pixely na okrajích (kde výška dosahuje limitu) reálně zahazuje podle vzoru, čímž udržuje stoprocentní čitelnost prostorové hloubky a správné vykreslování stínů.
\end{itemize}


\chapter{Výsledky a zhodnocení}
Tato kapitola se věnuje reálnému běhu vyvinuté aplikace, hodnocení jejího výkonu a celkové analýze zvolených programátorských přístupů. 
Hlavním kritériem úspěchu byla schopnost aplikace plynule vykreslovat a upravovat rozsáhlé matematické modely v prostředí virtuální reality.

\section{Testovací prostředí a metodika}
Pro objektivní zhodnocení výkonu probíhalo testování na hardwarové sestavě odpovídající střední třídě herních počítačů. 
Tato konfigurace poskytuje dostatečný výkon pro plynulý běh virtuální reality a zároveň umožňuje odhalit případné úzké hrdla neoptimalizovaného kódu.

Hardwarová specifikace testovací stanice:
\begin{itemize}
    \item \textbf{Grafická karta (GPU):} NVIDIA GeForce RTX 3060 Ti
    \item \textbf{Procesor (CPU):} AMD Ryzen 7 5700G
    \item \textbf{Operační paměť (RAM):} 16 GB
\end{itemize}

Zobrazení a interakce probíhaly primárně pomocí headsetu \product{Meta Quest 3}. 
Testování zahrnovalo obě běžně dostupné metody připojení k počítači. 
První metodou bylo přímé spojení kabelem zaručující maximální propustnost obrazu a minimální latenci. 
Druhou metodou byl bezdrátový přenos přes lokální síť (\term{Air Link}), který uživateli poskytuje naprostou volnost pohybu za cenu mírné komprese obrazových dat.

\section{Snímková frekvence a stabilita (FPS)}
Nejdůležitějším ukazatelem použitelnosti aplikace ve virtuální realitě je stabilita snímkové frekvence. V enginu \product{Unity} bylo provedeno měření pro různé hustoty mřížky. 
Sledovaly se dva odlišné stavy. 
Prvním byl klidový stav, kdy uživatel graf pouze pozoruje a prochází jím. 
Druhým stavem byla aktivní manipulace, kdy uživatel plynule táhne za posuvník parametru a aplikace musí v každém snímku generovat novou prostorovou síť.

\begin{table}[htbp]
    \centering
    \begin{tabular}{|c|c|c|c|}
        \hline
        \textbf{Rozlišení sítě} & \textbf{Počet vrcholů} & \textbf{FPS v klidu} & \textbf{FPS při změnách parametrů} \\
        \hline
        $100 \times 100$ & 10 000 & 72 & 55 \\
        $200 \times 200$ & 40 000 & 71 & 35 \\
        $250 \times 250$ & 62 500 & 70 & 30 \\
        $300 \times 300$ & 90 000 & 71 & 16 \\
        \hline
    \end{tabular}
    \caption{Závislost snímkové frekvence na hustotě sítě a uživatelské interakci.}
    \label{tab:vysledky_fps}
\end{table}

Z naměřených hodnot (viz tabulka \ref{tab:vysledky_fps}) vyplývá zásadní přínos architektury řízené událostmi. 
V klidovém stavu naráží aplikace na hardwarový limit obnovovací frekvence samotného headsetu (72 Hz). 
Tím se potvrzuje, že samotné vykreslení a zobrazení sítě s 90 000 vrcholy nedělá grafické kartě sebemenší problém. Výkon není limitován vykreslovacím vláknem.

Při aktivní manipulaci s parametry dochází k očekávanému snížení snímkové frekvence. 
Zde se projevuje výpočetní náročnost generování sítě na procesoru, zejména při iterování metodou bisekce. 
Zajímavým a klíčovým zjištěním však je, že uživatelský zážitek ve virtuální realitě zůstává plynulý i při naměřených hodnotách kolem 30 FPS. 

Nutno ale podotknout, že měření probíhalo přímo ve vývojovém prostředí \product{Unity Editor}, které samo o sobě spotřebovává značnou část výkonu. 
Ve finálním přeloženém sestavení aplikace se výpočetní limity posouvají o něco výše.


\section{Praktické ukázky vygenerovaných modelů}
Pro ověření funkčnosti celého řešení byla aplikace testována na sadě různorodých matematických předpisů. 
Cílem bylo potvrdit schopnost analyzátoru správně interpretovat složitý uživatelský textový řetězec a následně demonstrovat vizuální stabilitu topologicky odlišných ploch.

\subsection{Spojité plochy a dynamické parametry}
Prvním testovacím scénářem byla komplexní goniometrická funkce definovaná předpisem $f(x, y) = a \cdot \sin(b \cdot x) + c \cdot \cos(d \cdot y)$. 
Tento předpis primárně ověřuje schopnost aplikace vykreslovat a upravovat spojité funkce. 

Ve virtuální realitě vytváří tato rovnice plynulou zvlněnou plochu. 
Vykreslování a prostorová deformace sítě je zde bezproblémové i při velmi rychlé manipulaci s posuvníky (viz obrázek \ref{fig:ukazka_vlneni}). 
Funkce totiž na celém svém definičním oboru neobsahuje žádné body nespojitosti. 
Procesor tak není zatěžován náročnými iteracemi vyhledávací metody bisekce.

\begin{figure}[htbp]
    \centering
    \includegraphics[width=0.8\textwidth]{graf_sin_cos.png}
    \caption{Vizualizace spojité goniometrické funkce se čtyřmi modifikovatelnými parametry ($a, b, c, d$).}
    \label{fig:ukazka_vlneni}
\end{figure}

\subsection{Zpracování singularit a dělení nulou}
Druhým příkladem zastupujícím extrémně problematické rovnice je lomená funkce s předpisem $f(x, y) = \frac{a \cdot x^2 \cdot b \cdot y^2}{c \cdot x^3 + d \cdot y^2}$. 
Tento matematický model představuje zatěžkávací zkoušku pro samotný generátor geometrie.

Graf v tomto případě obsahuje výškové skoky a složité singularity v místech, kde se hodnota jmenovatele blíží k nule. 
Na výsledném vygenerovaném modelu (viz obrázek \ref{fig:ukazka_zlomek}) je jasně patrné fungování naprogramované topologické ochrany. 
Většinu nechtěných polygonů natahujících se přes propast generátor úspěšně detekuje a ořezává pomocí algoritmu bisekce. 
Při detailním zkoumání lze však v některých místech zpozorovat mírně zubaté okraje nebo ojedinělý polygon protínající oblast nespojitosti. 

Tato chyba je primárně způsobena konečným rozlišením generované sítě a maximálním povoleným počtem iterací vyhledávacího algoritmu. 
Pokud je skok funkce příliš strmý a reálná vzdálenost bodů napříč propastí je menší než základní krok mřížky, algoritmus nedokáže tuto hranu s absolutní přesností izolovat. 
Snaha o úplné odstranění tohoto jevu by vyžadovala drastické zahuštění sítě nebo neúměrné navýšení počtu iterací. 
To by následně vedlo ke kritickému propadu snímkové frekvence. 
Z kontextu celkového tvaru plochy je nicméně pro pozorovatele naprosto zřejmé, že tyto chybné spojnice do grafu matematicky nepatří a nepůsobí tak při průzkumu natolik rušivě. 

\begin{figure}[htbp]
    \centering
    \includegraphics[width=0.8\textwidth]{graf_zlomek.png}
    \caption{Ukázka lomené funkce.}
    \label{fig:ukazka_zlomek}
\end{figure}

\chapter{Závěr}
Cílem této práce bylo vyvinout aplikaci, která umožní vizualizovat a interaktivně zkoumat matematické funkce o dvou neznámých ve virtuální realitě. 
Aplikace by měla být schopna zpracovávat širokou škálu funkcí, včetně funkcí s body nespojitosti, a zároveň poskytovat plynulý uživatelský zážitek bez ohledu na komplexitu zobrazené geometrie.

\section{Shrnutí dosažených výsledků}
Aplikace úspěšně plní všechny stanovené cíle. 
Zcela bezproblémově zpracovává spojité funkce a s vysokou mírou spolehlivosti vizualizuje i funkce se singularitami. 
Tyto cíle byly dosaženy pomocí následujících přístupů a optimalizací:

\begin{itemize}
    \item \textbf{Algoritmus \term{Shunting-yard} pro analýzu rovnic:} Tento algoritmus umožnil efektivní převod uživatelského textového vstupu do formy, 
    která je snadno zpracovatelná pro vytváření abstraktního syntaktického stromu a jeho následné kompilace.
    \item \textbf{Kompilace pro rychlé vyhodnocování:} Vytvoření syntaktického stromu a jeho kompilace do mezikódu zrychluje výpočet jednotlivých vrcholů prostorové sítě na procesoru.
    \item \textbf{Algoritmus bisekce:} Implementace algoritmu bisekce zajišťuje detekci a ořezávání polygonů protínajících propasti, což umožňuje bezpečné zobrazování funkcí s body nespojitosti.
    \item \textbf{Optimalizace paměti a recyklace objektů:} Minimalizace alokací a využití existujících datových struktur urychluje překreslování grafu a výrazně snižuje zátěž procesu \term{Garbage Collector}.
    \item \textbf{Architektura řízená událostmi:} Oddělení výpočtů od renderovací smyčky zajišťuje, že procesorový čas je využíván pouze v okamžicích změny dat. 
    Díky tomu si aplikace udržuje maximální snímkovou frekvenci v klidovém stavu.
\end{itemize}

\section{Možnosti budoucího rozšíření}
Ačkoliv aplikace úspěšně plní svůj hlavní účel, nabízí se několik možností pro její budoucí rozšíření a optimalizaci:

\begin{itemize}
    \item \textbf{Komplexnější zpracování singularit:} Ačkoliv aplikace aktuálně dokáže detekovat a odstranit většinu chybných polygonů, v místech nejsložitějších nespojitostí se stále objevují drobné vizuální artefakty.
    \item \textbf{Nezávislé sestavení:} V současné verzi je aplikace určena pro běh primárně přes počítač. 
    Překlad do čistě nezávislého sestavení by naplno využil potenciál headsetů s vlastním výpočetním výkonem (jako je například \product{Meta Quest}).
    \item \textbf{Rozhraní pro zadávání funkcí:} Místo externí úpravy textového souboru by bylo uživatelsky přívětivější implementovat virtuální klávesnici přímo do aplikace, 
    což by umožnilo zadávat a ukládat nové rovnice bez nutnosti sundávat headset.
\end{itemize}
\appendix


\chapter{Uživatelská příručka}
Tato příloha slouží jako návod k ovládání aplikace v prostředí virtuální reality. 
Uživatelské rozhraní je navrženo tak, aby bylo co nejvíce intuitivní. 
Uživatel k pohybu a manipulaci s objekty využívá standardní pohybové ovladače podporované vývojovým rámcem \product{XR Interaction Toolkit}.

\begin{figure}[htbp]
    \centering
    \includegraphics[width=0.7\textwidth]{ovladac.png}
    \caption{Základní rozložení prvků a interakčních paprsků na virtuálním ovladači.}
    \label{fig:manual_ovladac}
\end{figure}

\section{Hlavní tabule a výběr funkce}
Po spuštění aplikace se uživatel ocitne ve virtuálním prostoru, kde na stěně se nachází hlavní interaktivní tabule. 
Na této ploše se nachází seznam všech dostupných matematických funkcí. 
Tyto funkce se dynamicky načítají z externího textového souboru. 

\begin{figure}[htbp]
    \centering
    \includegraphics[width=0.9\textwidth]{tabule.png}
    \caption{Hlavní interaktivní tabule s nabídkou načtených matematických funkcí.}
    \label{fig:manual_tabule}
\end{figure}

Pro výběr konkrétní rovnice slouží virtuální laserový paprsek (\term{Raycast}). 
Tento paprsek (viz obrázek \ref{fig:manual_laser}) vychází nepřetržitě z virtuálních pohybových ovladačů. 
Uživatel jednoduše zamíří na požadovanou rovnici a stiskne \texttt{Trigger}. 
Okamžitě poté proběhne na pozadí výpočet a v prostoru se vygeneruje trojrozměrný model zvoleného grafu.

\begin{figure}[htbp]
    \centering
    \includegraphics[width=0.6\textwidth]{laser.png}
    \caption{Ukázka virtuálního laserového paprsku sloužícího k výběru funkce.}
    \label{fig:manual_laser}
\end{figure}

\section{Úprava seznamu funkcí}
Nabídka rovnic na hlavní tabuli není v aplikaci pevně zakódována. 
Veškeré matematické předpisy se dynamicky načítají z externího textového souboru. 
Tento přístup umožňuje uživatelům velmi snadno přidávat, upravovat nebo mazat vlastní matematické modely bez nutnosti jakkoliv zasahovat do zdrojového kódu.

\begin{figure}[htbp]
    \centering
    \includegraphics[width=0.7\textwidth]{ukazka_textaku.png}
    \caption{Ukázka struktury textového souboru s definicemi matematických funkcí.}
    \label{fig:manual_soubor}
\end{figure}

Proces přidání nové funkce do aplikace je navržen s ohledem na maximální plynulost práce a skládá se z následujících kroků:

\begin{itemize}
    \item \textbf{Lokalizace a otevření souboru:} Uživatel vyhledá příslušný textový soubor v kořenové složce aplikace a otevře jej v libovolném textovém editoru (například v programu Poznámkový blok).
    \item \textbf{Zápis rovnice:} Každá matematická funkce musí být zapsána výhradně na samostatném řádku. 
    Aplikace plně podporuje standardní matematický zápis včetně závorek a běžných operací (například \command"sin", \command"cos", \command"sqrt"). 
    Vstupní prostorové proměnné jsou striktně označovány písmeny $x$ a $y$.
    \item \textbf{Využití vlastních parametrů:} Pokud uživatel do zápisu vloží libovolné neznámé znaky (například písmena $a, b, c$), analyzátor je při načítání automaticky rozpozná jako dynamické parametry. 
    Pro tyto nalezené proměnné se posléze ve virtuální realitě samy vygenerují posuvníky v příručním menu.
    \item \textbf{Uložení a automatická aktualizace:} Jakmile uživatel textový soubor uloží, aplikace tuto změnu na pozadí okamžitě detekuje. Seznam na virtuální tabuli se automaticky překreslí.
\end{itemize}


\section{Fyzická manipulace s modelem}
Vygenerovaný matematický objekt představuje plně interaktivní prvek. 
Aplikace naplno využívá fyzikální povahu virtuální reality a umožňuje s grafem prostorově manipulovat následovně:

\begin{itemize}
    \item \textbf{Přesun a rotace:} Uživatel může ke grafu přistoupit a stisknutím úchopového tlačítka (viz obrázek \ref{fig:manual_uchop}) jej vzít přímo do ruky. 
    Graf lze posléze libovolně otáčet, naklánět a přesouvat v prostoru pro detailní prozkoumání ze všech úhlů.
    
    \begin{figure}[htbp]
        \centering
        \includegraphics[width=0.5\textwidth]{tlacitko_uchop.png}
        \caption{Detail úchopového tlačítka na virtuálním ovladači.}
        \label{fig:manual_uchop}
    \end{figure}

    \item \textbf{Změna měřítka:} Pro zkoumání asymptot nebo naopak celkového tvaru funkce je k dispozici režim obouručního škálování. Uživatel uchopí graf oběma rukama současně. 
    Následným plynulým pohybem rukou od sebe se model zvětšuje. Pohybem rukou k sobě se model naopak zmenšuje.
    
    \item \textbf{Obnova pozice:} Může nastat situace, kdy uživatel nedopatřením upustí graf mimo vymezenou herní plochu, případně jej příliš zmenší. 
    Pro tyto případy je k dispozici resetovací tlačítko (viz obrázek \ref{fig:manual_reset}), které pozici a měřítko modelu okamžitě obnoví.
    
    \begin{figure}[htbp]
        \centering
        \includegraphics[width=0.5\textwidth]{tlacitko_reset.png}
        \caption{Tlačítko pro okamžitý reset pozice a měřítka grafu.}
        \label{fig:manual_reset}
    \end{figure}
\end{itemize}

\section{Příruční menu a dynamické úpravy}
Pro hlubší analýzu tvaru a chování funkce má uživatel k dispozici dvě interaktivní příruční menu. 
Toto rozhraní je ukotveno k virtuálnímu levému ovladači a vyvolává se stisknutím určitého tlačítka. Každé z těchto menu slouží k jiné úpravě grafu.

\begin{figure}[htbp]
    \centering
    \includegraphics[width=0.9\textwidth]{menu_parametry.png}
    \caption{Příruční menu ovladače pro změnu parametrů}
    \label{fig:manual_param_menu}
\end{figure}
\begin{figure}[htbp]
    \centering
    \includegraphics[width=0.9\textwidth]{menu_def.png}
    \caption{Příruční menu ovladače pro změnu definičního oboru.}
    \label{fig:manual_def_menu}
\end{figure}

První menu slouží k úpravě parametrů rovnice a uživatel jej vyvolá tlačítkem ukázaným na obrázku \ref{fig:manual_btn_parametry}. 
Nacházejí se zde dynamicky vygenerované posuvníky pro úpravu nalezených koeficientů (například parametrů $a, b, c$). 
Uživatel je může laserovým paprskem uchopit a plynule s nimi pohybovat. Každá změna hodnoty se okamžitě projeví na grafu.

\begin{figure}[htbp]
    \centering
    \includegraphics[width=0.5\textwidth]{tlacitko_parametry.png}
    \caption{Tlačítko pro vyvolání menu s parametry funkce.}
    \label{fig:manual_btn_parametry}
\end{figure}

Druhé menu upravuje nastavení definičního oboru a aktivuje se příslušným tlačítkem (viz obrázek \ref{fig:manual_btn_obor}). 
Obsahuje krokovací tlačítka pro přesné určení hranic vykreslování na osách $x$ a $y$. 
\term{Uniform} pak umožňuje měnit velikost definičního oboru symetricky do všech směrů současně.

\begin{figure}[htbp]
    \centering
    \includegraphics[width=0.5\textwidth]{tlacitko_obor.png}
    \caption{Tlačítko pro vyvolání menu s nastavením definičního oboru.}
    \label{fig:manual_btn_obor}
\end{figure}

\section{Průzkum vnitřní topologie}
Zásadní výhodou virtuální reality je volnost pohybu. 
Uživatel není vázán pouze na statický pohled zvenčí. 
Díky implementaci oboustranného grafického materiálu může do vygenerovaného modelu fyzicky vstoupit. 
Získá tak zcela unikátní perspektivu pro pochopení složitých prostorových útvarů a vnitřních struktur zkoumaných matematických funkcí.





\backmatter
\printbibliography
\backpage
\end{document}